
\subsection{Green 张量统一前向模型}\label{subsec:green-forward-closure-dp47}

线性介质对任意电流源的响应已经由推迟 Green 张量确定。自由电子实验只是改变了源和读出方式：电子自身作为源时，可以统计它失去的总能量，也可以只统计传到远场的辐射；外加相干光场作为源时，则读取电子获得的相干能量调制。因此三类观测共享同一个频域响应核，只在源项和探测投影上不同。

\begin{definition}[电子源与散射响应]
对给定经典电子轨迹 $\vb{r}_e(t)$，电流 Fourier 分量为
\begin{equation}
 \widetilde{\vb{J}}_e(\vb{r},\omega)
 =q_{\mathrm e}\int\!\dd t\,
 \vb{v}(t)\delta^{(3)}[\vb{r}-\vb{r}_e(t)]\ee^{\ii\omega t}.
\label{eq:electron_current_fourier_dp47}
\end{equation}
由线性响应产生的诱导场为
\begin{equation}
{
 \widetilde{\vb{E}}_{\rm ind}(\vb{r},\omega)
 =\ii\mu_0\omega
 \int\!\dd^3 r'\,
 \mathbf G_{\rm sc}^R(\vb{r},\vb{r}';\omega)
 \widetilde{\vb{J}}_e(\vb{r}',\omega),}
\label{eq:induced_field_green_dp47}
\end{equation}
其中 $\mathbf G_{\rm sc}^R=\mathbf G^R-\mathbf G_0^R$ 只保留环境相对于参考真空的散射响应。

\end{definition}

\begin{proposition}[轨迹上的损失泛函]
电子单位角频率损失概率可由电子对诱导场所做功得到。这里的虚部按算符虚部理解，核为 $[\mathbf G^R(\vb{r},\vb{r}')-\mathbf G^R(\vb{r}',\vb{r})^\dagger]/(2\ii)$；在本章互易介质中等于逐分量虚部。对匀速直线轨迹 $\vb{r}_e(t)=\vb{R}_0+\vb{v}_0t$，写成 Green 张量双时间积分为
\begin{equation}
{
 \Gamma_{\rm EELS}(\omega)
 =\frac{\mu_0q_{\mathrm e}^2}{\pi\hbar}
 \int\!\dd t\,\dd t'\,
 \ee^{\ii\omega(t-t')}
 \vb{v}_0\!\cdot
 \operatorname{Im}\mathbf G_{\rm sc}^R
 (\vb{r}_e(t),\vb{r}_e(t');\omega)
 \cdot\vb{v}_0 .}
\label{eq:eels_green_master_dp47}
\end{equation}
$\Gamma_{\rm EELS}(\omega)\dd\omega$ 是无量纲概率，因此 $[\Gamma_{\rm EELS}]={\rm s}$。

\end{proposition}

\begin{definition}[收集通道振幅与光子谱]
CL 不是把\eqnrefs{eq:eels_green_master_dp47} 中的 $\operatorname{Im}\mathbf G^R$ 换成模方，而是先把电子电流经 Green 张量传播到远场探测通道。令探测模式 $\lambda$ 的归一化接收函数为 $\vb{f}_\lambda(\vb{r},\omega)$，定义通道振幅
\begin{equation}
 \mathcal A_\lambda(\omega)
 =\ii\mu_0\omega
 \int\!\dd^3 r\,\dd^3 r'\,
 \vb{f}_\lambda^*(\vb{r},\omega)\!\cdot
 \mathbf G^R(\vb{r},\vb{r}';\omega)
 \widetilde{\vb{J}}_e(\vb{r}',\omega).
\label{eq:cl_channel_amplitude_dp47}
\end{equation}
将接收函数归一化为谱能量振幅，使 $|\mathcal A_\lambda|^2$ 的单位为 $\mathrm{J\,s}$、表示每单位角频率的通道能量后，光子谱密度为
\begin{equation}
{
 \Gamma_{\rm CL}^{\rm det}(\omega)
 =\sum_{\lambda\in\mathcal C}
 \frac{|\mathcal A_\lambda(\omega)|^2}{\hbar\omega},}
\label{eq:cl_detected_green_master_dp47}
\end{equation}
其中集合 $\mathcal C$ 编码数值孔径、偏振、频率窗和空间模式匹配。这样“总辐射”与“被探测 CL”严格区分。

\end{definition}

\begin{proposition}[相干驱动的轨迹积分]
受激 PINEM 则由外加相干场沿电子轨迹所做的功决定。对沿 $z$ 方向以速度 $v_0$ 传播的电子，时域功为
\begin{equation}
{
 \Delta E_{\rm ext}
 =q_{\mathrm e}\int_{-\infty}^{\infty}\!\dd t\,
 v_0 E_{{\rm ext},z}(\vb{R}_0,v_0t;t).}
\label{eq:pinem_time_domain_work_dp68}
\end{equation}
若外场只含角频率 $\omega$，定义正频幅 $\vb{E}_{\rm ext}^{(+)}(\vb{r})$，使实场为 $\vb{E}_{\rm ext}^{(+)}\ee^{-\ii\omega t}+\mathrm{c.c.}$。它的单位是 ${\rm V/m}$，不同于脉冲的连续 Fourier 谱幅 ${\rm V\,s/m}$。沿轨迹的经典无量纲耦合为
\begin{equation}
{
 \beta(\omega)
 =\frac{q_{\mathrm e}}{\hbar\omega}
 \int_{-\infty}^{\infty}\!\dd z\,
 E_{{\rm ext},z}^{(+)}(\vb{R}_0,z)
 \ee^{-\ii\omega z/v_0}.}
\label{eq:pinem_beta_green_unified_dp47}
\end{equation}
若外场由具有相同正频幅约定的单频电流 $\vb{J}_{\rm ext}^{(+)}$ 经同一环境产生，则将
\begin{equation*}
 \vb{E}_{\rm ext}^{(+)}
 =\ii\mu_0\omega\,\mathbf G^R*\vb{J}_{\rm ext}^{(+)}
\end{equation*}
代入\eqnrefs{eq:pinem_beta_green_unified_dp47}，PINEM 也成为同一个 $\mathbf G^R$ 的线性泛函。若改用实场 $\operatorname{Re}[\vb{E}_{\rm ph}\ee^{-\ii\omega t}]$ 的通常相量，则 $\vb{E}_{\rm ext}^{(+)}=\vb{E}_{\rm ph}/2$；不能把连续 Fourier 电流谱直接当作该单频幅度。

\end{proposition}

\begin{remark}[同一响应核的三种读取]
因此三类实验的结构可压缩为
\begin{equation}
{
 \mathbf G^R
 \longrightarrow
 \begin{cases}
 \operatorname{Im}\mathbf G^R & \text{EELS：电子不可逆损失},\\
 \Pi_{\mathcal C}\mathbf G^R\widetilde{\vb{J}}_e & \text{CL：辐射到收集通道},\\
 \mathbf G^R\vb{J}_{\rm ext}^{(+)} & \text{PINEM：外场相干驱动},
 \end{cases}}
\label{eq:green_three_experiments_map_dp47}
\end{equation}
这里的差别来自输入源和最终读取的观测量，而不是三种不同的介质理论。

\end{remark}

\begin{definition}[仪器响应后的计数似然]
实际反演还要把理论谱映射到探测器计数。把待估参数记为 $\vb{\vartheta}$，理论连续谱记为 $s(\omega|\vb{\vartheta})$，仪器响应矩阵为 $R_{j}(\omega)$，则第 $j$ 个像素的期望计数为
\begin{equation}
{
 \lambda_j(\vb{\vartheta})
 =N_e\int\!\dd\omega\,
 R_j(\omega)s(\omega|\vb{\vartheta})+b_j,}
\label{eq:detector_count_forward_dp47}
\end{equation}
其中 $N_e$ 是入射电子数，$b_j$ 是背景期望计数。若入射采用 Poisson 模型并经独立分流，像素计数独立，完整似然为
\begin{equation}
{
 \log\mathcal L(\vb{\vartheta})
 =\sum_j\left[n_j\log\lambda_j-\lambda_j-\log(n_j!)\right].}
\label{eq:poisson_likelihood_green_inverse_dp47}
\end{equation}
这使 Green 张量、电子态、收集模式、仪器分辨率与最终参数反演处在同一个前向模型中。

\end{definition}

\begin{proof}[证明（匀速损失泛函的约化）]
一般的“做功 $\to\operatorname{Im}\mathbf G^R\to$ 损失概率”关系已经由\eqnrefs{eq:prescribed_current_green_loss_master,eq:green_tensor_eels_master} 建立。对具体几何，只需把相应响应核代入一般表达式：使用散射 Green 张量 $\mathbf G_{\rm sc}^R=\mathbf G^R-\mathbf G_0^R$ 去除参考真空贡献，并取匀速轨迹
\[
\vb{r}_e(t)=\vb{R}_0+\vb{v}_0t,
\qquad
\vb{J}_e(\vb{r},t)=q_{\mathrm e}\vb{v}_0\delta^{(3)}[\vb{r}-\vb{r}_e(t)].
\]
一般表达式随即化成\eqnrefs{eq:eels_green_master_dp47}。由此可以直接看到：材料和几何只改变 Green 前向核，电子跃迁理论仍由同一个电流--响应双线性控制。
\end{proof}
